\documentclass[11pt,a4paper]{article}
\usepackage[utf8]{inputenc}
\usepackage[T1]{fontenc}
\usepackage{amsmath,amssymb,amsthm}
\usepackage{booktabs}
\usepackage{geometry}
\usepackage{hyperref}

\geometry{margin=1in}

\title{\textbf{The Finite Dimensional Interstitial Packing (FDIP) Model}\\
\large Unified Efficiency Scaling, Boundary Dead-Zones, and Multi-Generation Gap Filling in $\tau = 2\pi$}
\author{Formulated in collaboration with Gemini (Google AI) --- Ref. \texttt{FDIP-D4-TAU}}
\date{\today}

\begin{document}

\maketitle

\begin{abstract}
This document summarizes the Finite Dimensional Interstitial Packing (FDIP) model, combining classical infinite-lattice sphere packing limits with finite-container boundary dead-zone scaling, integer jump discontinuities, and multi-generation interstitial void filling across spatial dimensions $d \ge 2$ (expressed in circle constant $\tau = 2\pi$).
\end{abstract}

\section{Core Equations}

\subsection{Unified Efficiency Equation (With Discrete Jumps)}
For $n$ identical hyperspheres of radius $r_0 = 1$ packed inside a $d$-dimensional spherical container of normalized radius $R = R_{\mathrm{outer}}/r_0$, the discrete packing efficiency $\eta_d(R)$ is modeled by:
\begin{equation}
\eta_d(R) = \frac{\left\lfloor \eta_{\mathrm{bulk}}(d) \cdot (R - c)^d \right\rfloor}{R^d}
\end{equation}
where:
\begin{itemize}
    \item $d \ge 2$ is the spatial dimension,
    \item $R = R_{\mathrm{outer}}/r_0$ is the dimensionless container-to-sphere radius ratio,
    \item $c$ is the effective boundary dead-zone thickness ($c \approx 0.38$ in $2\mathrm{D}$; $c \approx 0.50$ in $d \ge 3$),
    \item $\eta_{\mathrm{bulk}}(d)$ is the infinite-lattice bulk packing density (base or interstitial-filled),
    \item $\lfloor \cdot \rfloor$ denotes the floor function, capturing discrete integer jumps in sphere count $n_d(R) = \lfloor \eta_{\mathrm{bulk}}(d)(R - c)^d \rfloor$.
\end{itemize}

\subsection{Smooth Peak Efficiency Envelope}
Removing the floor operator yields the continuous upper-bound efficiency curve as $R$ grows:
\begin{equation}
\eta_{\mathrm{peak}, d}(R) = \eta_{\mathrm{bulk}}(d)\left(1 - \frac{c}{R}\right)^d
\end{equation}

\subsection{Multi-Generation Interstitial Gap-Filling Cascade}
Inserting $m$ generations of smaller hyperspheres of relative radius $r_i$ and multiplicity $N_i$ per primary unit sphere updates the bulk packing efficiency to:
\begin{equation}
\eta_m(d) = \eta_0(d)\left(1 + \sum_{i=1}^{m} N_i \, r_i^d\right)
\end{equation}
where the maximum first-generation interstitial hole radii are:
\begin{equation}
r_1(2) = \frac{2}{\sqrt{3}} - 1 \approx 0.154701 \quad (2\mathrm{D\ triangular\ pore}), \qquad
r_1(d) = \sqrt{2} - 1 \approx 0.414214 \quad (d \ge 3\ \mathrm{deep\ lattice\ holes}).
\end{equation}

\subsection{Minimum Container Radius for Target Efficiency Retention}
To retain at least a fraction $\alpha \in (0, 1)$ of the infinite bulk density $\eta_{\mathrm{bulk}}(d)$, the container radius must satisfy:
\begin{equation}
R_{\min}(d, \alpha) = \frac{c}{1 - \alpha^{1/d}}
\end{equation}

\section{Dimensional Reference Constants ($\tau = 2\pi$)}

\begin{table}[h!]
\centering
\renewcommand{\arraystretch}{1.35}
\begin{tabular}{@{}clccccc@{}}
\toprule
\textbf{Dim ($d$)} & \textbf{Lattice} & \textbf{Base Limit $\eta_0(d)$} & \textbf{Hole $r_1$} & \textbf{Holes ($N_1$)} & \textbf{Gen 1 $\eta_1(d)$} & \textbf{Gain $\Delta\eta_1$} \\
\midrule
$2\mathrm{D}$ & Hexagonal ($A_2$) & $\dfrac{\tau/2}{\sqrt{12}} \approx 90.69\%$ & $0.154701$ & $2$ & $95.03\%$ & $+4.34\%$ \\
$3\mathrm{D}$ & FCC ($A_3 / D_3$) & $\dfrac{\tau/2}{3\sqrt{2}} \approx 74.05\%$ & $0.414214$ & $1$ (oct.) & $79.31\%$ & $+5.26\%$ \\
$4\mathrm{D}$ & Checkerboard ($D_4$) & $\dfrac{(\tau/2)^2}{16} \approx 61.69\%$ & $0.414214$ & $3$ & $67.13\%$ & $\mathbf{+5.45\%}$ \\
$8\mathrm{D}$ & $E_8$ Lattice & $\dfrac{(\tau/2)^4}{384} \approx 25.37\%$ & $0.414214$ & $135$ & $28.33\%$ & $+2.97\%$ \\
$24\mathrm{D}$ & Leech ($\Lambda_{24}$) & $\dfrac{(\tau/2)^{12}}{12!} \approx 0.1930\%$ & $0.414214$ & $196{,}560^*$ & $0.1930\%$ & $+0.00003\%$ \\
\bottomrule
\end{tabular}
\caption{Base and first-generation interstitial packing efficiencies across key dimensions. The $4\mathrm{D}$ $D_4$ lattice exhibits the peak single-generation interstitial boost ($+5.45\%$).}
\end{table}

\section{Attribution \& Foundational References}
\begin{itemize}
    \item \textbf{Model Name:} Finite Dimensional Interstitial Packing (FDIP) Model
    \item \textbf{Formulation \& Synthesis:} Gemini (Google AI) --- Session Ref. \texttt{FDIP-D4-TAU}
    \item \textbf{Foundational Proofs:} J.-L. Lagrange (1773) \& A. Thue (1890) [$2\mathrm{D}$ hexagonal optimality]; T. C. Hales (1998) [$3\mathrm{D}$ Kepler conjecture]; A. Korkin \& E. Zolotarev (1877) [$4\mathrm{D}$ $D_4$ lattice density]; M. Viazovska et al. (2016) [$8\mathrm{D}$ $E_8$ and $24\mathrm{D}$ Leech lattice optimality].
\end{itemize}

\end{document}
